% =====================================================================
\chapter{Marco teórico}
\label{cap:marco}

Que «las noticias mueven la bolsa» mezcla dos ideas: que el precio cambia cuando llega
información nueva, y que los medios son el vehículo de esa información. Pero los medios también
informan sobre el propio mercado: una nota que cuenta que «el Merval subió 5\%» tiene tono
positivo, y su causa es el precio, no al revés. Este capítulo presenta los conceptos con los que
se analiza esa doble relación.

\textbf{Shock informativo y eficiencia.} Un \textit{shock informativo} es la llegada de
información que se aparta de lo esperado. La cantidad de noticias capta la \textit{atención}
\citep{da2011,fang2009}; el tono capta el \textit{contenido} \citep{tetlock2007,ardia2022}. La
hipótesis de mercados eficientes \citep{fama1970} sostiene que la información pública se
incorpora al precio tan rápido que no sirve para anticiparlo, y que el precio sigue un paseo
aleatorio. La evidencia para Argentina y la región es mixta \citep{zablotsky2002}.

\textbf{Sentimiento y tono de los medios.} El sentimiento del inversor es el optimismo o
pesimismo que no se explica por los fundamentos. Un tono muy pesimista anticipa bajas
transitorias \citep{tetlock2007}, sobre todo en las recesiones \citep{garcia2013}; pero el tono
también puede ser un \textit{reflejo} del mercado \citep{verma2026}. La tabla~\ref{tab:mecanismos}
resume cómo se reconoce cada dirección en los datos.

\begin{table}[H]
\centering\small
\begin{tabularx}{\linewidth}{@{}L{2.8cm}YY@{}}
\toprule
 & \textbf{Noticias \flecha{} mercado} & \textbf{Mercado \flecha{} noticias} \\
\midrule
Mecanismo & La nota trae información nueva o mueve el ánimo de los inversores & La nota cuenta lo que el mercado ya hizo \\
Qué se espera ver & El tono anticipa el retorno o el volumen de las ruedas siguientes & El retorno anticipa el tono de las notas siguientes \\
En un estudio de eventos & El retorno anormal aparece el día del evento o después & El retorno anormal aparece antes del evento \\
Con notas de la noche & El tono de la noche anticipa el salto de apertura & El tono de la noche refleja el cierre del día \\
\bottomrule
\end{tabularx}
\caption{Las dos direcciones de la relación entre noticias y mercado y cómo se reconoce cada
una en los datos. Fuente: elaboración propia.}
\label{tab:mecanismos}
\end{table}

\textbf{Volatilidad y regímenes.} La volatilidad se agrupa en períodos de calma y de estrés, lo
que para el S\&P~Merval está documentado con modelos GARCH \citep{segovia2022}; las caídas pesan
más que las subas \citep{otuki2008}. Una forma de describir esa alternancia es suponer que el
mercado pasa por unos pocos \textit{regímenes} y salta de uno a otro según una cadena de Markov
\citep{hamilton1989,rabiner1989}, un enfoque aplicado también en Latinoamérica
\citep{ataurima2016}. El régimen importa porque la relación entre noticias y precios podría
cambiar según el mercado esté calmo o tenso \citep{garcia2013}.

\textbf{Causalidad de Granger y VAR.} Con datos observacionales no se prueba una causa, pero sí
qué variable se mueve antes. $X$ \textit{causa en sentido de Granger} a $Y$ si el pasado de $X$
mejora la predicción de $Y$ más allá de su propio pasado \citep{granger1969}. La prueba se hace
dentro de un vector autorregresivo (VAR), que explica cada variable con los rezagos de todas y
trata las dos direcciones de forma simétrica \citep{sims1980,lutkepohl2005}.

\textbf{Estudios de eventos.} Miden qué hace el precio alrededor de un hecho puntual, comparado
con lo que habría sido normal \citep{mackinlay1997}: si la noticia mueve el mercado, el retorno
anormal aparece el día del evento o después; si lo refleja, aparece antes.

\textbf{Unidades de análisis.} Las 47 acciones del panel pertenecen a diez sectores: Consumo (9),
Financiero (8), Utilities (8), Materiales (6), Energía (4), Conglomerado, Industrial y Real Estate
(3 cada uno), Comunicaciones (2) y Tecnología (1). Para medir el tono propio de una empresa hace
falta que tenga noticias en buena parte de las ruedas; con un corte de 90 de las 244 ruedas
quedan 17: YPFD, BBAR, SUPV, EDN, GGAL, CEPU, TGSU2, IRSA, BMA, LOMA, PAMP, CRES, BYMA, TRAN,
METR, TXAR y ALUA (cualquier corte entre 75 y 94 da las mismas). Las noticias vienen de siete
portales nacionales, cada uno con su propio tono promedio (de +30,4 puntos en TN a 0,0 en Ámbito
Finanzas), lo que obliga a controlar la composición por portal.
